\section{Thiết kế hệ thống}

\subsection{Kiến trúc tổng quan}

SmartLock Fuzzy được thiết kế theo mô hình backend trung tâm. Backend sở hữu camera và phần cứng, còn frontend chỉ đóng vai trò giao diện quan sát/điều khiển qua HTTP API với FastAPI và Next.js \cite{fastapi,nextjs}.


\begin{figure}[H]
\centering
\resizebox{\textwidth}{!}{%
\begin{tikzpicture}[
    >=Stealth,
    font=\small,
    block/.style={
        draw,
        rounded corners,
        minimum width=3.3cm,
        minimum height=0.9cm,
        align=center
    },
    hw/.style={
        draw,
        rounded corners,
        minimum width=3.0cm,
        minimum height=0.85cm,
        align=center
    },
    data/.style={thick, -{Stealth}},
    http/.style={thick, {Stealth}-{Stealth}},
    ctrl/.style={thick, dashed, -{Stealth}},
    lab/.style={
        font=\scriptsize,
        fill=white,
        inner sep=1.5pt
    }
]

% ---------------- Nodes ----------------
\node[hw]    (cam)          at (0,  1.4) {Camera\\\texttt{/dev/video0}};
\node[hw]    (hardware)     at (0, -1.4) {Keypad/OLED/Servo\\GPIO + SPI + PWM};

\node[block] (worker)       at (4.2, 0) {\texttt{LocalCameraWorker}\\vòng lặp xử lý};

\node[block] (api)          at (8.5, 1.4) {FastAPI\\REST endpoints};
\node[block] (front)        at (12.7, 1.4) {Next.js Dashboard\\polling HTTP};

\node[block] (vision)       at (8.5, -0.4) {OpenCV\\Haar + LBPH};
\node[block] (fuzzy)        at (8.5, -2.2) {Fuzzy Controller\\Mamdani rules};

\node[block] (registration) at (4.2, -2.2) {Registration Manager\\thu mẫu + train LBPH};

% ---------------- Data flow ----------------
\draw[data] (cam.east) -- node[lab, above] {frame} (worker.west);

\draw[data] (worker.north east) -- node[lab, above] {latest frame/status} (api.west);
\draw[http] (api.east) -- node[lab, above] {GET/POST} (front.west);

\draw[data] (worker.east) -- node[lab, above] {frame} (vision.west);
\draw[data] (vision.south) -- node[lab, right] {detection} (fuzzy.north);

\draw[data] (worker.south) -- node[lab, right] {samples} (registration.north);

% ---------------- Hardware control ----------------
\draw[ctrl] (worker.south west) to[bend right=15]
    node[lab, left] {unlock/status} (hardware.east);

\draw[ctrl] (hardware.east) to[bend right=15]
    node[lab, below] {keypad event} (worker.west);

\end{tikzpicture}
}
\caption{Kiến trúc tổng quan của hệ thống SmartLock Fuzzy}
\label{fig:system_architecture}
\end{figure}

Các thành phần chính:

\begin{itemize}
    \item \textbf{Backend FastAPI}: quản lý vòng đời hệ thống, API và worker camera.
    \item \textbf{Camera worker}: xử lý frame, nhận diện, đánh giá fuzzy và cập nhật trạng thái mới nhất.
    \item \textbf{Service xử lý ảnh}: đóng gói Haar Cascade, LBPH và các phép đo chất lượng khuôn mặt.
    \item \textbf{Service fuzzy}: ánh xạ kết quả nhận diện thành quyết định truy cập.
    \item \textbf{Service phần cứng}: điều khiển keypad, OLED và servo.
    \item \textbf{Frontend}: dashboard web để quan sát camera, đăng ký người dùng và đổi mật khẩu.
\end{itemize}

\subsection{Luồng xử lý camera}

Mỗi frame đi qua các bước sau:

\begin{enumerate}
    \item Đọc frame từ camera bằng OpenCV \cite{opencv}.
    \item Lật ảnh nếu biến môi trường \texttt{CAMERA\_FLIP} được bật.
    \item Phát hiện khuôn mặt bằng Haar Cascade \cite{viola2001}.
    \item Nếu có nhiều khuôn mặt, giữ khuôn mặt có bounding box lớn nhất.
    \item Nhận diện bằng LBPH nếu model đã được nạp \cite{lbph2006}.
    \item Tính độ sáng ROI và góc lệch tương đối.
    \item Đưa detection đầu tiên vào bộ fuzzy để tính \texttt{security\_risk}.
    \item Nếu hành động là \texttt{unlock} và keypad không hoạt động, kích hoạt servo mở khóa.
    \item Xử lý đăng ký khuôn mặt nếu đang có session đăng ký.
    \item Vẽ overlay, mã hóa JPEG base64 và cập nhật cache API.
\end{enumerate}

\begin{figure}[H]
\centering
\resizebox{\textwidth}{!}{%
\begin{tikzpicture}[
    >=Stealth,
    node distance=0.85cm,
    step/.style={draw, rounded corners, minimum width=3.0cm, minimum height=0.8cm, align=center, font=\small},
    arrow/.style={thick, -{Stealth}}
]
    \node[step] (read) {Đọc frame};
    \node[step, right=of read] (detect) {Haar detect};
    \node[step, right=of detect] (lbph) {LBPH recognize};
    \node[step, right=of lbph] (quality) {Đo sáng/góc};
    \node[step, below=of quality] (fuzzy) {Fuzzy risk};
    \node[step, left=of fuzzy] (action) {Quyết định};
    \node[step, left=of action] (encode) {Overlay + JPEG};
    \node[step, left=of encode] (api) {API cache};

    \draw[arrow] (read) -- (detect);
    \draw[arrow] (detect) -- (lbph);
    \draw[arrow] (lbph) -- (quality);
    \draw[arrow] (quality) -- (fuzzy);
    \draw[arrow] (fuzzy) -- (action);
    \draw[arrow] (action) -- (encode);
    \draw[arrow] (encode) -- (api);
\end{tikzpicture}
}
\caption{Pipeline xử lý một frame camera}
\label{fig:camera_pipeline}
\end{figure}

\subsection{Thiết kế API}

Backend cung cấp các endpoint chính sau:

\begin{table}[H]
\centering
\caption{Danh sách endpoint của hệ thống}
\label{tab:api_endpoints}
\begin{tabularx}{\textwidth}{|l|l|X|}
\hline
\textbf{Method} & \textbf{Endpoint} & \textbf{Chức năng} \\
\hline
GET & \texttt{/health} & Kiểm tra backend còn hoạt động. \\
\hline
GET & \texttt{/camera/status} & Trạng thái camera worker, đăng ký và phần cứng. \\
\hline
GET & \texttt{/api/v1/camera/frame} & Metadata và frame JPEG base64 mới nhất. \\
\hline
GET & \texttt{/api/v1/camera/latest} & Metadata mới nhất, không kèm frame base64. \\
\hline
GET & \texttt{/api/v1/camera/history} & Lịch sử kết quả gần nhất. \\
\hline
POST & \texttt{/api/v1/register/start} & Bắt đầu thu mẫu đăng ký khuôn mặt. \\
\hline
POST & \texttt{/api/v1/register/cancel} & Hủy session đăng ký hiện tại. \\
\hline
GET & \texttt{/api/v1/register/status} & Lấy tiến trình đăng ký. \\
\hline
POST & \texttt{/api/v1/keypad/password} & Đổi mã PIN 6 chữ số. \\
\hline
GET & \texttt{/api/v1/keypad/status} & Lấy trạng thái keypad và lockout. \\
\hline
\end{tabularx}
\end{table}

Cache camera mới nhất được lưu trong module router và bảo vệ bằng \texttt{threading.Lock}. Cách này đơn giản, đủ dùng cho một backend cục bộ phục vụ dashboard.

\subsection{Thiết kế fuzzy decision}

Bộ fuzzy nhận detection đầu tiên. Nếu không có detection, hệ thống trả về \texttt{deny}. Nếu danh tính là \texttt{Unknown}, độ tin cậy mô hình được đặt bằng 0. Với người nhận diện được, khoảng cách LBPH được scale thành \texttt{model\_confidence}. Bộ fuzzy dùng 13 luật theo hướng Mamdani \cite{mamdani1975}:

\begin{itemize}
    \item LOW confidence luôn dẫn tới MAXIMUM risk.
    \item HIGH confidence, ánh sáng NORMAL và mặt FRONTAL dẫn tới MINIMUM risk.
    \item HIGH confidence nhưng ánh sáng/góc không tối ưu dẫn tới AVERAGE risk.
    \item MEDIUM confidence chỉ được xem là AVERAGE risk trong một số điều kiện thuận lợi; các điều kiện xấu tăng lên MAXIMUM risk.
\end{itemize}

Khi thư viện fuzzy không khả dụng, \texttt{services/fuzzy\_logic.py} có fallback logic tĩnh để backend không bị dừng. Fallback vẫn giữ nguyên tinh thần fail-safe: thiếu tin cậy thì rủi ro cao.

\subsection{Thiết kế đăng ký khuôn mặt}

Đăng ký khuôn mặt là một session có trạng thái. Người dùng nhập tên trên dashboard, backend tạo thư mục tương ứng trong \texttt{registered\_faces/}. Mỗi frame được kiểm tra chất lượng trước khi lưu:

\begin{itemize}
    \item Chỉ chấp nhận đúng một khuôn mặt.
    \item Kích thước khuôn mặt đủ lớn.
    \item Độ sáng trong khoảng an toàn.
    \item Góc mặt không vượt ngưỡng.
    \item Ảnh không quá mờ.
    \item Mẫu mới phải khác mẫu trước đó đủ lớn để tránh lưu các frame trùng lặp.
\end{itemize}

Khi đủ số mẫu, backend huấn luyện lại LBPH, ghi file XML và JSON label map, sau đó reload recognizer mà không cần khởi động lại server.

\subsection{Thiết kế phần cứng}

Bảng sau mô tả mapping GPIO mặc định:

\begin{table}[H]
\centering
\caption{Kết nối phần cứng Raspberry Pi}
\label{tab:hardware_pins}
\begin{tabular}{|l|l|l|}
\hline
\textbf{Thiết bị} & \textbf{Chức năng} & \textbf{GPIO BCM} \\
\hline
Keypad 4x4 & Row 1--4 & 17, 27, 22, 5 \\
\hline
Keypad 4x4 & Col 1--4 & 6, 13, 19, 26 \\
\hline
Servo & PWM signal & 18 \\
\hline
OLED SSD1306 SPI & DC, RST & 24, 25 \\
\hline
OLED SSD1306 SPI & MOSI, SCLK, CE0 & 10, 11, 8 \\
\hline
\end{tabular}
\end{table}

\begin{figure}[H]
\centering
\resizebox{\textwidth}{!}{%
\begin{tikzpicture}[
    >=Stealth,
    font=\small,
    pi/.style={draw, rounded corners, fill=green!8, minimum width=4.0cm, minimum height=5.8cm, align=center},
    dev/.style={draw, rounded corners, fill=blue!7, minimum width=3.6cm, minimum height=1.0cm, align=center},
    power/.style={draw, rounded corners, fill=red!7, minimum width=3.6cm, minimum height=1.0cm, align=center},
    wire/.style={thick, -{Stealth}},
    pwr/.style={thick, red!70!black, -{Stealth}},
    gnd/.style={thick, black!75, -{Stealth}},
    sig/.style={thick, blue!70!black, -{Stealth}},
    note/.style={font=\scriptsize, align=center, fill=white, inner sep=1pt}
]

\node[pi] (pi) at (0,0) {\textbf{Raspberry Pi 4B}\\[0.2cm]
\begin{tabular}{ll}
GPIO17 & Keypad R1\\
GPIO27 & Keypad R2\\
GPIO22 & Keypad R3\\
GPIO5  & Keypad R4\\
GPIO6  & Keypad C1\\
GPIO13 & Keypad C2\\
GPIO19 & Keypad C3\\
GPIO26 & Keypad C4\\
GPIO18 & Servo PWM\\
GPIO10 & OLED MOSI\\
GPIO11 & OLED SCLK\\
GPIO8  & OLED CS\\
GPIO24 & OLED DC\\
GPIO25 & OLED RST
\end{tabular}};

\node[dev] (cam) at (-5.6,2.4) {USB Camera\\\texttt{/dev/video0}};
\node[dev] (keypad) at (5.8,2.2) {Keypad 4x4\\R1--R4, C1--C4};
\node[dev] (oled) at (5.8,0.0) {OLED SSD1306 SPI\\VCC, GND, SCLK, MOSI, RST, DC, CS};
\node[dev] (servo) at (5.8,-2.2) {RC Servo\\Signal, 5V, GND};
\node[power] (supply) at (-5.6,-2.2) {Nguồn 5V ngoài\\cho servo};

\draw[wire] (cam.east) -- node[note, above] {USB} (pi.west);
\draw[sig] (pi.east) -- node[note, above] {8 GPIO matrix} (keypad.west);
\draw[sig] (pi.east) -- node[note, above] {SPI + DC/RST} (oled.west);
\draw[sig] (pi.east) -- node[note, above] {GPIO18 PWM} (servo.west);
\draw[pwr] (supply.east) -- node[note, below] {5V} (servo.west);
\draw[gnd] (supply.north east) to[bend left=18] node[note, above] {GND chung} (pi.south west);
\draw[gnd] (pi.south east) to[bend right=18] node[note, below] {GND} (servo.south west);

\end{tikzpicture}
}
\caption{Sơ đồ kết nối phần cứng của hệ thống SmartLock Fuzzy}
\label{fig:hardware_schematic}
\end{figure}

Keypad sử dụng interrupt trên các cột để tránh polling liên tục. Servo được điều khiển bằng PWM, sau đó duty cycle được đưa về 0 để giảm rung. OLED được render bằng Pillow trước khi gửi qua SPI qua thư viện luma.oled \cite{lumaoled}.

\subsection{Thiết kế frontend}

Frontend Next.js gồm ba thành phần chính:

\begin{itemize}
    \item \texttt{Header}: tiêu đề và trạng thái dashboard.
    \item \texttt{LiveVideoStream}: polling frame, hiển thị camera, thông tin nhận diện, fuzzy risk và form đăng ký.
    \item \texttt{PasswordChange}: đổi mật khẩu keypad, kiểm tra độ dài 6 chữ số và gửi API.
\end{itemize}

Frontend đọc biến môi trường \texttt{NEXT\_PUBLIC\_API\_URL}, mặc định là \texttt{http://localhost:8000}. Dashboard polling frame mỗi 250 ms, phù hợp cho demo live view cục bộ.
